\section{Conclusion}

Spurious correlations degrade worst-group accuracy while maintaining high average accuracy, yet existing detection methods require knowing which spurious feature to search for \cite{adebayo2022post}. Robust learning approaches either demand group annotations (GroupDRO, JTT) or operate post-hoc on embeddings (SCER). Gradient attribution methods fail for unknown spurious correlations --- practitioners cannot detect spurious reliance using saliency maps alone.

We proposed a gradient-based framework that detects and mitigates spurious correlations via gradient abnormality, extending GAIA \cite{chen2023exploring} from distribution shift (OOD) to subpopulation shift (minority groups). Our methodology operates on gradient \emph{process} (abnormality, not attribution) and provides unified detection-mitigation in a single pipeline.

\textbf{Validated contributions:}

\textbf{Detection (h-e1):} Gradient abnormality pipeline (GradCAM $\rightarrow$ GAIA-Z $\rightarrow$ statistical testing) differentiates minority vs majority patterns in synthetic experiments (divergence 0.30, $p<0.0001$, Cohen's $d=198.75$). First application of GAIA to spurious correlation detection.

\textbf{Mechanism (h-m-integrated):} Causal validation via synthetic tests shows (1) strong correlation between GAIA divergence and WGA ($|\rho|=0.975$, $p<0.001$), (2) background augmentation reduces GAIA-Z by 39.4\% ($p<0.001$, $d=2.96$), supporting spurious-conflict hypothesis.

\textbf{Mitigation (h-m-mitigate):} Spatial gradient regularization improves WGA by +23pp on MNIST+Color toy dataset (78\% vs 55\%) in proof-of-concept experiments without catastrophic accuracy drop.

\textbf{Limitations acknowledged:} All validation from synthetic data (h-e1, h-m-integrated) or minimal PoC (h-m-mitigate). Real Waterbirds experiments require GPU compatibility resolution (PyTorch 2.1+ for H100 sm\_90 support) or CPU training ($\sim$70-90 GPU hours). Empirical claims (real minority gradients exhibit abnormality, real Waterbirds WGA improvement) pending full-scale validation.

\textbf{Contribution tier:} Tier 3 (methodology validated) with clear path to Tier 2 (real empirical validation). Synthetic validation proves code correctness and mechanism plausibility; real validation required for empirical claim verification.

\textbf{Immediate future work:} (1) Real Waterbirds detection/mitigation experiments, (2) Full 5-seed MNIST validation, (3) GroupDRO baseline comparison. \textbf{Extensions:} Multi-dataset generalization (CelebA, UrbanCars), joint gradient-embedding regularization, automatic hyperparameter selection.

\textbf{Callback to hook:} We showed gradient abnormality \emph{can} differentiate minority patterns when the spurious-conflict mechanism holds (synthetic validation). Next step: Does real Waterbirds data exhibit this mechanism? (Pending infrastructure upgrade and full-scale experiments.)

The gradient space encodes training dynamics --- what models rely on during learning. Our work demonstrates gradient abnormality as a viable paradigm for spurious correlation detection, complementing embedding-based approaches and bypassing attribution ineffectiveness. Methodology validated; empirical validation in progress.
